\documentclass[11pt]{article}

\usepackage[T1]{fontenc}
\usepackage{amsmath,amssymb,amsthm}
\usepackage[margin=1in]{geometry}
\usepackage[hidelinks]{hyperref}
\setlength{\emergencystretch}{2em}

\newtheorem{theorem}{Theorem}
\newtheorem{lemma}[theorem]{Lemma}
\newtheorem{corollary}[theorem]{Corollary}
\newtheorem{conjecture}[theorem]{Conjecture}
\theoremstyle{remark}
\newtheorem{remark}[theorem]{Remark}

\newcommand{\R}{\mathbb R}
\newcommand{\D}{\mathcal D}
\newcommand{\Plo}[1]{P_{\le #1}}

\title{A Carrier-Turnover Clock for Type-II Navier--Stokes Records}
\author{John Robert Lawson and Codex}
\date{Research round ending 24 July 2026}

\begin{document}
\maketitle

\begin{abstract}
This round proves an energy-only low-frequency temporal estimate for smooth
three-dimensional Navier--Stokes flow and applies it to the active Type-II
programme.  If energy-normalised carrier profiles form a nonzero compact
subset of \(L^2(\R^3)\), replacing a radius-\(R\) carrier by a sufficiently
smaller carrier requires time at least \(cR^{5/2}\), uniformly in moving
centres.  The exact representative \(q=4\) ledger has
\(R_j\asymp2^{-4j}\) but record gaps
\(\asymp2^{-11j}/j=o(R_j^{5/2})\).  It therefore has no infinite nonzero
strongly \(L^2\)-precompact carrier-profile subsequence on one NSE trajectory.
The result is conditional
on that representative schedule and compactness; it forces the surviving
branch into nested microstructure, spatial escape, fragmentation, or another
loss of compactness.  It does not solve the Clay problem.
\end{abstract}

\section{Scope and change from the preceding round}

The preceding round constructed, for each small scale separately, a smooth
compact NSE packet with fixed Gaussian-band work on its \(r^{5/2}\) turnover
time and only \(O(\sqrt r)\) dissipation.  It showed that the local
half-radius work charge is sharp, but did not place successive packets on
one trajectory.

This report tests the most literal concatenation: a single trajectory whose
energy-normalised carrier states remain in a compact nonzero family while
their radii shrink on the exact surviving record schedule.  That
configuration is impossible.  The proof has two ingredients:

\begin{enumerate}
\item strict \(L^2\)-critical dilation evacuates every fixed low-frequency
window, uniformly on compact profile families and uniformly in translations;
\item the projected NSE equation bounds the speed of any low-frequency
change using only the kinetic-energy ceiling.
\end{enumerate}

All claims below are theorem-level relative to their displayed assumptions
but await external mathematical review.  They are new within this
repository; no literature-novelty claim is made.

\section{Uniform low-frequency separation}

For \(R>0\) and \(z\in\R^3\), define
\[
 (\D_{R,z}f)(x)
 :=
 R^{-3/2}f\!\left(\frac{x-z}{R}\right).
\]
This action is unitary on \(L^2(\R^3)\).

Fix a radial function \(\chi\in C_c^\infty(\R^3)\) satisfying
\[
 0\le\chi\le1,\qquad
 \chi=1\ \hbox{on }B_1,\qquad
 \operatorname{supp}\chi\subset B_2,
\]
and let \(\Plo{\Lambda}\) be the Fourier multiplier
\(\chi(\xi/\Lambda)\).

\begin{lemma}[Compact profiles separate from strict dilations]
\label{lem:separation}
Let \(\mathcal K\subset L^2(\R^3)\) be relatively compact and suppose
\[
 d:=\inf_{f\in\mathcal K}\|f\|_2>0.
\]
There are \(K<\infty\), \(\lambda_0\in(0,1)\), and \(\delta>0\) such that
for all \(f,g\in\mathcal K\), all \(z\in\R^3\), and all
\(0<\lambda\le\lambda_0\),
\[
 \left\|
 \Plo K\bigl(f-\D_{\lambda,z}g\bigr)
 \right\|_2
 \ge\delta.
\]
\end{lemma}

\begin{proof}
The strong convergence \(\Plo K\to I\) as \(K\to\infty\) is uniform on
compact subsets of \(L^2\).  Choose \(K\) so that
\[
 \sup_{f\in\mathcal K}
 \|(I-\Plo K)f\|_2\le\frac d8.
\]
Then
\[
 \inf_{f\in\mathcal K}\|\Plo Kf\|_2\ge\frac{7d}{8}.
\]

With a unitary Fourier transform,
\[
 \widehat{\D_{\lambda,z}g}(\xi)
 =
 \lambda^{3/2}e^{-iz\cdot\xi}\widehat g(\lambda\xi).
\]
Changing variables gives
\[
\begin{aligned}
 \|\Plo K\D_{\lambda,z}g\|_2^2
 &=
 \int_{\R^3}
 \left|\chi\!\left(\frac{\eta}{\lambda K}\right)\right|^2
 |\widehat g(\eta)|^2\,d\eta\\
 &\le
 \int_{|\eta|\le2\lambda K}|\widehat g(\eta)|^2\,d\eta.
\end{aligned}
\]
The right-hand side is independent of \(z\), tends to zero for each \(g\),
and does so uniformly on \(\mathcal K\).  For the last assertion, cover the
compact closure by a finite \(L^2\) net and use the uniform operator bound
\(\|\Plo K\D_{\lambda,z}\|_{2\to2}\le1\).  Thus
\[
 \sup_{g\in\overline{\mathcal K},\,z\in\R^3}
 \|\Plo K\D_{\lambda,z}g\|_2
 \longrightarrow0.
\]
Choose \(\lambda_0\) so that this supremum is at most \(d/8\).  The triangle
inequality then gives the assertion with \(\delta=3d/4\).
\end{proof}

\begin{remark}
The uniformity in \(z\) is important.  A carrier cannot evade the conclusion
merely by moving its spatial centre between records.
\end{remark}

\section{An energy-only NSE low-pass clock}

Let \(u\) be a smooth divergence-free solution on \([s,t]\) of
\[
 \partial_tu+\mathbb P\nabla\!\cdot(u\otimes u)
 =
 \nu\Delta u
\]
on \(\R^3\), and write
\[
 E_*:=\sup_{\tau\in[s,t]}\frac12\|u(\tau)\|_2^2.
\]

\begin{theorem}[Low-frequency temporal Lipschitz bound]
\label{thm:lowpass}
For every \(\Lambda>0\),
\[
 \boxed{
 \|\Plo{\Lambda}(u(t)-u(s))\|_2
 \le
 C_\chi
 \left(
 E_*\Lambda^{5/2}
 +
 \nu\sqrt{E_*}\Lambda^2
 \right)(t-s).
 }
\]
\end{theorem}

\begin{proof}
Integrating the Leray-projected equation,
\[
\begin{aligned}
 \Plo{\Lambda}(u(t)-u(s))
 ={}&
 -\int_s^t
 \Plo{\Lambda}\mathbb P\nabla\!\cdot(u\otimes u)(\tau)\,d\tau\\
 &+
 \nu\int_s^t\Plo{\Lambda}\Delta u(\tau)\,d\tau.
\end{aligned}
\]
Each component of the nonlinear operator has multiplier
\[
 \chi(\xi/\Lambda)i\xi_\ell
 \left(\delta_{im}-\frac{\xi_i\xi_m}{|\xi|^2}\right).
\]
By Plancherel its convolution kernel has \(L^2\) norm
\(C_\chi\Lambda^{5/2}\).  Young's convolution inequality therefore gives
\[
\begin{aligned}
 \|\Plo{\Lambda}\mathbb P\nabla\!\cdot(u\otimes u)\|_2
 &\le
 C_\chi\Lambda^{5/2}\|u\otimes u\|_1\\
 &\le
 C_\chi\Lambda^{5/2}\|u\|_2^2
 \le
 2C_\chi E_*\Lambda^{5/2}.
\end{aligned}
\]
The viscous multiplier has magnitude at most \(C_\chi\Lambda^2\), whence
\[
 \nu\|\Plo{\Lambda}\Delta u\|_2
 \le
 C_\chi\nu\Lambda^2\|u\|_2
 \le
 C_\chi\nu\sqrt{2E_*}\Lambda^2.
\]
Integration in time and enlargement of \(C_\chi\) prove the estimate.
\end{proof}

\begin{corollary}[Carrier-turnover lower bound]
\label{cor:clock}
If \(R>0\) and
\[
 \|\Plo{K/R}(u(t)-u(s))\|_2\ge\delta,
\]
then
\[
 t-s
 \ge
 \frac{\delta R^{5/2}}
 {C_\chi\left(
 E_*K^{5/2}
 +
 \nu\sqrt{E_*}K^2R^{1/2}
 \right)}.
\]
For fixed data and \(R\downarrow0\), this is at least \(cR^{5/2}\).
\end{corollary}

\begin{proof}
Set \(\Lambda=K/R\) in Theorem~\ref{thm:lowpass} and rearrange.
\end{proof}

\section{The compact-carrier theorem}

Let \(t_j\uparrow T^*\), \(R_j\downarrow0\), and let \(x_j\in\R^3\) be
arbitrary.  Define the exact energy-normalised carrier profiles
\[
 F_j(y)
 :=
 R_j^{3/2}u(x_j+R_jy,t_j).
\]
Equivalently,
\[
 u(\cdot,t_j)=\D_{R_j,x_j}F_j,
 \qquad
 \|F_j\|_2=\|u(t_j)\|_2.
\]

\begin{theorem}[Compact carriers cannot outrun turnover]
\label{thm:carrier}
Suppose \(\{F_j\}\) is relatively compact in \(L^2(\R^3)\) and
\[
 \inf_j\|F_j\|_2>0.
\]
There exist \(\lambda_0,c>0\) such that, for all sufficiently large
\(j<k\) satisfying \(R_k/R_j\le\lambda_0\),
\[
 \boxed{t_k-t_j\ge cR_j^{5/2}.}
\]
The constants are independent of \(x_j\) and \(x_k\).
\end{theorem}

\begin{proof}
Apply Lemma~\ref{lem:separation} to the closure of \(\{F_j\}\).  In the old
carrier coordinates,
\[
 \D_{R_j,x_j}^{-1}\D_{R_k,x_k}F_k
 =
 \D_{\lambda,z}F_k,
 \quad
 \lambda=\frac{R_k}{R_j},
 \quad
 z=\frac{x_k-x_j}{R_j}.
\]
Low-pass covariance and unitarity yield
\[
\begin{aligned}
 &\left\|
 \Plo{K/R_j}\bigl(u(t_j)-u(t_k)\bigr)
 \right\|_2\\
 &\qquad=
 \left\|
 \Plo K\bigl(F_j-\D_{\lambda,z}F_k\bigr)
 \right\|_2
 \ge\delta.
\end{aligned}
\]
Corollary~\ref{cor:clock} gives
\[
 t_k-t_j
 \ge
 \frac{\delta R_j^{5/2}}
 {C_\chi\left(
 E_*K^{5/2}
 +
 \nu\sqrt{E_*}K^2R_j^{1/2}
 \right)}.
\]
The energy ceiling is fixed and \(R_j\to0\), so the denominator is uniformly
bounded for large \(j\).
\end{proof}

\section{The exact \texorpdfstring{\(q=4\)}{q=4} representative}

The existing joint power ledger takes, up to fixed comparison constants,
\[
 R_j=2^{-4j},
 \qquad
 t_{j+1}-t_j=\frac{c_t2^{-11j}}j.
\]

\begin{corollary}[No compact-profile subsequence for the \(q=4\) ledger]
\label{cor:q4}
No smooth finite-energy NSE trajectory can obey this scale and time schedule
while an infinite subsequence of its profiles \(F_j\) is relatively compact
in \(L^2\) and bounded away from zero.
\end{corollary}

\begin{proof}
Suppose such a subsequence exists and apply Theorem~\ref{thm:carrier} to it.
For every sufficiently large retained index \(j\), choose a later retained
\(k\) such that
\[
 \frac{R_k}{R_j}\le\lambda_0.
\]
Theorem~\ref{thm:carrier} requires
\[
 t_k-t_j\ge cR_j^{5/2}=c2^{-10j}.
\]
But
\[
 t_k-t_j
 \le
 T^*-t_j
 =
 \sum_{n\ge j}\frac{c_t2^{-11n}}n
 \le
 C\frac{2^{-11j}}j,
\]
so
\[
 \frac{t_k-t_j}{R_j^{5/2}}
 \le
 C\frac{2^{-j}}j
 \longrightarrow0.
\]
This is a contradiction.  The proof is unchanged when the displayed powers
hold with fixed two-sided comparison constants.
\end{proof}

In the carrier notation already used in the dossier,
\[
 v_j(y,0)
 =
 \frac1{a_j}u(x_j+R_jy,t_j),
 \qquad
 b_j=a_j^2R_j^3,
\]
and therefore
\[
 F_j=\sqrt{b_j}\,v_j(\cdot,0).
\]
The amplitude layer in an energy-efficient cell bounds \(\|F_j\|_2\) away
from zero.  Thus even subsequential strong \(L^2\) compactness of the carrier
states, after the usual scalar subsequence for \(b_j\), is excluded by
Corollary~\ref{cor:q4}.

\section{Why the subgrid branch survives}

The conclusion concerns compactness at the carrier radius \(R_j\).  The
actual defect band in the representative ledger has
\[
 r_j=R_j\ell_j\asymp2^{-5j}.
\]
Its fixed-energy local turnover time is
\[
 r_j^{5/2}\asymp2^{-25j/2},
\]
whereas the record gap is \(2^{-11j}/j\).  Hence
\[
 \frac{2^{-11j}/j}{r_j^{5/2}}
 =
 \frac{2^{3j/2}}j
 \longrightarrow\infty.
\]
A sufficiently small packet already present below the carrier scale has
ample local turnover time.  The theorem therefore does not contradict the
preceding sharp packet construction.  It forces any fast survivor to hide
the next event in noncompact subgrid structure rather than generating a
fresh compact carrier after the preceding event.

\section{Findings, limitations, and open obligations}

\subsection*{Robust findings, subject to review}

\begin{enumerate}
\item Bounded kinetic energy gives the unconditional low-pass temporal bound
of Theorem~\ref{thm:lowpass}.
\item Compact nonzero \(L^2\) carrier families have a uniform
translation-independent low-frequency mismatch after sufficient strict
dilation.
\item Combining the two facts forces the \(R^{5/2}\) carrier clock.
\item The exact \(q=4\) representative outruns that clock by
\(2^{-j}/j\), so it has no infinite nonzero strongly precompact
carrier-profile subsequence.
\end{enumerate}

\subsection*{Things still to prove}

\begin{enumerate}
\item Derive a usable profile decomposition for the forced failure of
\(L^2\) compactness.
\item Separate frequency ancestry from spatial escape and fragmentation.
\item Prove a nonsummable charge, rigidity theorem, or energy-vanishing
alternative for every resulting defect branch.
\item Treat schedules other than the representative \(q=4\) powers.
\item Keep the divergent-normalised-energy carrier cells separate.
\end{enumerate}

\begin{conjecture}[Nested-ancestry charge]
For an energy-efficient first-record Type-II sequence satisfying
\[
 t_{j+N}-t_j=o(R_j^{5/2})
\]
across a fixed nontrivial scale contraction, the noncompact part of the
carrier profile decomposition contains either a spatial-import event with a
fixed physical cost or a nested high-frequency descendant.  Along an
infinite descendant chain, those costs cannot all be summable against the
finite energy-dissipation budget.
\end{conjecture}

The first sentence of this conjecture is motivated by the theorem but is not
proved by it; the nonsummability assertion is the major open step.

\section{Epistemic conclusion}

This round closes one concrete subbranch: naive compact-profile
concatenation cannot realise the exact fast survivor.  It does not exclude a
noncompact same-trajectory cascade, produce a singularity, prove regularity
or breakdown, or establish any Clay alternative A--D.  The Clay problem
remains unsolved.

\end{document}
